%%=============================================================================
%% Conclusie
%%=============================================================================

\chapter{Conclusie}%
\label{ch:conclusie}

% Trek een duidelijke conclusie, in de vorm van een antwoord op de
% onderzoeksvra(a)g(en). Wat was jouw bijdrage aan het onderzoeksdomein en
% hoe biedt dit meerwaarde aan het vakgebied/doelgroep? 
% Reflecteer kritisch over het resultaat. In Engelse teksten wordt deze sectie
% ``Discussion'' genoemd. Had je deze uitkomst verwacht? Zijn er zaken die nog
% niet duidelijk zijn?
% Heeft het onderzoek geleid tot nieuwe vragen die uitnodigen tot verder 
%onderzoek?

%TODO: aanvullen en nakijken
Dit onderzoek vertrok vanuit een alledaags maar fundamenteel probleem: 
de discrepantie tussen de werkelijke operationele status van een supermarkt 
en de weergave ervan in digitale POI-databanken. De centrale onderzoeksvraag 
luidde: \textit{"Hoe kan artificiële intelligentie worden ingezet om automatisch 
    de (historische) open of gesloten status van Points-of-Interest te detecteren, 
    en hoe kan deze aanpak zorgen voor actuelere data in locatiegebaseerde 
    toepassingen?"}. In dit hoofdstuk worden de bevindingen per deelvraag 
samengevat, gevolgd door een algemene conclusie en aanbevelingen voor 
toekomstig onderzoek.

\section{Antwoord op de deelvragen}

\subsection{Welke gegevensbronnen kunnen worden gebruikt?}

Zoals besproken in Hoofdstuk~\ref{ch:stand-van-zaken} vereist een robuuste 
aanpak een combinatie van complementaire bronnen: structurele POI-attributen 
als feitelijke basis, dynamische mobiliteitsbronnen voor vroegtijdige 
signaaldetectie, straatbeeldbronnen voor visuele verificatie en crowd-sourced 
platformen als aanvullende annotatie. In de proof-of-concept werd gewerkt 
met de interne POI-databank van Accurat, aangevuld met historische 
statusregistraties en gestructureerde openingstijden. Het ontbreken van 
voetgangers- of transactiegegevens bleef een structurele beperking van 
de huidige feature-set.

\subsection{Wat zijn de uitdagingen bij de implementatie van AI?}

De voornaamste uitdagingen, uitvoerig besproken in 
Hoofdstuk~\ref{ch:stand-van-zaken}, komen ook in de praktijk duidelijk 
naar voren. De datakwaliteit en -actualiteit vormen de meest fundamentele 
hindernis: records zijn frequent gefragmenteerd, inconsistent of verouderd. 
Data-inflatie bemoeilijkt het labelproces en verhoogt de kans op foutieve 
trainingsdata. Het cold-start probleem maakt nauwkeurige classificatie 
van jonge vestigingen structureel moeilijker, omdat gedragsfeatures zoals 
de \texttt{inactiviteit\_ratio} voor nieuwe POI's onbetrouwbaar zijn. 
Ten slotte vormt geospatiale bias een reëel risico: de dataset is sterk 
geconcentreerd op Franse supermarktketens, wat de generaliseerbaarheid 
naar kleinere regio's beperkt.

\subsection{Welke AI-technieken zijn geschikt?}

Op basis van de literatuurstudie en de proof-of-concept kan worden gesteld 
dat supervised classificatiemodellen een praktisch inzetbare aanpak vormen. 
In tegenstelling tot time series-benaderingen en computervision-methoden 
zijn tabellaire classificatiemodellen direct toepasbaar op intern beschikbare 
POI-data, zonder nood aan externe dynamische bronnen. De vergelijkende 
evaluatie van Logistische Regressie, Random Forest en XGBoost toont aan 
dat [TODO: kernbevinding op basis van de resultaten, bv. 
\textit{``ensemble-methoden significant beter presteren dan het lineaire 
    basismodel''}]. De feature-importantieanalyse bevestigt dat historisch 
gedragsmatige kenmerken — en niet uitsluitend temporele attributen — 
de sterkste predictoren zijn voor de gesloten klasse.

\subsection{Hoe kan de betrouwbaarheid worden gevalideerd?}

De betrouwbaarheidsvalidatie vereist een meerlagige aanpak. Op modelmatig 
niveau werden precision, recall en F1-score voor de klasse \textit{gesloten} 
als primaire evaluatiecriteria gehanteerd, zoals toegelicht in 
Hoofdstuk~\ref{ch:proof-of-concept}. Een ideale validatie maakt gebruik 
van een onafhankelijke ground truth-dataset, die in dit onderzoek 
ontbrak. Voor productieomgevingen is een continue monitoringlaag 
aangewezen die modeldrift detecteert en periodieke hertraining initieert.

\section{Algemene conclusie}

Deze bachelorproef toont aan dat het automatisch detecteren van de 
operationele status van supermarkten via supervised machine learning 
technisch haalbaar is op basis van intern beschikbare POI-data. De 
ontwikkelde proof-of-concept levert een werkend classificatiesysteem 
dat de status van een POI met [TODO: accuracy X\%] kan voorspellen.

Tegelijkertijd maakt dit onderzoek duidelijk dat de kwaliteit van de 
onderliggende data bepalend is voor de prestaties van elk 
statusdetectiesysteem. Zolang historische registraties onvolledig zijn, 
data-inflatie optreedt en dynamische mobiliteitsdata ontbreken, blijft 
de detectiecapaciteit structureel beperkt. De POI-levenscyclus — van 
\textit{booming} over \textit{stable} tot \textit{decaying} — kan worden 
weerspiegeld in gedragsfeatures, maar vereist een rijkere en continue 
datastroom om volledig te worden gevat.

\section{Aanbevelingen voor toekomstig onderzoek}

Op basis van de bevindingen worden vier concrete aanbevelingen geformuleerd. 
Ten eerste verdient de integratie van dynamische mobiliteitsbronnen 
prioriteit, aangezien voetgangers- of transactiedata de vroegste 
signalen van een statuswijziging bevatten. Ten tweede biedt de toepassing 
van deep learning-architecturen zoals LSTM-netwerken perspectief voor 
het verwerken van temporele statussequenties. Ten derde dient geospatiale 
bias in vervolgonderzoek expliciet te worden geadresseerd door 
steekproefstratificatie op land- en merkniveau toe te passen. Ten slotte 
vormt de evolutie van een statisch batchmodel naar een continu, 
zelflerend systeem — waarbij het model proactief statuswijzigingen 
detecteert en valideert — de meest waardevolle langetermijnrichting 
voor dit onderzoeksdomein.

